%%%%%%%%%%%%%%%%%%%%%%%%%%%%%%%%%%%%%%%%%%%%%%%%%%%%%%%%%%%%%%%%%%%%%%
% Lớp 4 | Chủ đề 1. Ôn tập và bổ sung | Ôn tập các số đến 100 000
% Dạng 4D1?1-1: Cấu tạo số, giá trị chữ số theo hàng
%%%%%%%%%%%%%%%%%%%%%%%%%%%%%%%%%%%%%%%%%%%%%%%%%%%%%%%%%%%%%%%%%%%%%%


% nguồn: 2026-11-18_lop4_chu-de-1-tu-luan (nhập 18/11/2026)
\begin{bt}%[4D1N1-1]%[ma:904aa7d3]
% nguon: Toán tư duy lớp 4 – Bài 1 SGK Toán 4 KNTT
	Số $70\,435$ gồm mấy chục nghìn, mấy nghìn, mấy trăm, mấy chục và mấy đơn vị? Viết số đó thành tổng theo các hàng.
	\loigiai{
		Đọc từ trái sang phải: chữ số $7$ ở hàng chục nghìn, $0$ ở hàng nghìn, $4$ ở hàng trăm, $3$ ở hàng chục, $5$ ở hàng đơn vị.\\
		Vậy số $70\,435$ gồm $7$ chục nghìn, $0$ nghìn, $4$ trăm, $3$ chục và $5$ đơn vị.\\
		$70\,435 = 70\,000 + 400 + 30 + 5$.\\
		Mẹo: hàng nào có chữ số $0$ thì không cần viết vào tổng.
	}
\end{bt}

% nguồn: 2026-11-18_lop4_chu-de-1-tu-luan (nhập 18/11/2026)
\begin{bt}%[4D1H1-1]%[ma:301d7f40]
% nguon: Toán tư duy lớp 4 – Bài 1 SGK Toán 4 KNTT
	Trong số $58\,358$, giá trị của chữ số $5$ ở hàng chục nghìn gấp bao nhiêu lần giá trị của chữ số $5$ ở hàng chục?
	\loigiai{
		Chữ số $5$ ở hàng chục nghìn có giá trị $50\,000$.\\
		Chữ số $5$ ở hàng chục có giá trị $50$.\\
		$50\,000 : 50 = 1\,000$.\\
		Vậy giá trị đó gấp $1\,000$ lần.
	}
\end{bt}

% nguồn: 2026-11-18_lop4_chu-de-1-tu-luan (nhập 18/11/2026)
\begin{bt}%[4D1V1-1]%[ma:1e98b1b5]
% nguon: Toán tư duy lớp 4 – Bài 1 SGK Toán 4 KNTT
	Tìm số lớn nhất có năm chữ số khác nhau mà tổng các chữ số của nó bằng $20$.
	\loigiai{
		Muốn số lớn nhất thì chữ số ở hàng cao phải lớn nhất có thể.\\
		Hàng chục nghìn chọn $9$, hàng nghìn chọn $8$. Tổng còn lại cho ba chữ số cuối: $20 - 9 - 8 = 3$.\\
		Ba chữ số khác nhau (khác $9$ và $8$) có tổng bằng $3$ chỉ có thể là $2$, $1$, $0$. Xếp giảm dần: $2$, $1$, $0$.\\
		Vậy số cần tìm là $98\,210$ (thử lại: $9+8+2+1+0=20$).
	}
\end{bt}

% nguồn: 2026-11-18_lop4_chu-de-1-tu-luan-bo-sung (nhập 18/11/2026)
\begin{bt}%[4D1H1-1]%[ma:de98eeaf]
% nguon: Toán tư duy lớp 4 – Bài 1 SGK Toán 4 KNTT
	Nếu thay chữ số $5$ của số $45\,678$ bằng chữ số $9$ thì được số mới. Số mới hơn số cũ bao nhiêu đơn vị?
	\loigiai{
		Chữ số $5$ ở hàng nghìn, có giá trị $5\,000$. Thay bằng $9$ thì có giá trị $9\,000$.\\
		Số mới là $49\,678$, hơn số cũ: $9\,000 - 5\,000 = 4\,000$ (đơn vị).\\
		Thử lại: $49\,678 - 45\,678 = 4\,000$.
	}
\end{bt}
